\section{The complete quartic directional layer of Tornheim zeta}
\label{qt:sec}

This section answers the explicit quartic-layer question in the incoming
\emph{Cubic Tornheim Derivatives and Harmonic Stieltjes Constants},
Section 7, Question 4 \cite{RepoCubic}, and the corresponding question in
\emph{Independent Orders and Exact Reductions}, Section 7 \cite{RepoIndependent}.
The incoming work already proves the cubic ray formula and evaluates its
diagonal coordinate in a convergent harmonic series. Those results are
not new here. The present extension classifies the next homogeneous layer,
gives convergent representatives for every coordinate, and obtains exact
quartic cancellation identities.

Put
\[
 L=\log(2\pi),\qquad d_j=\zeta^{(j)}(0),\qquad
 e_j=\zeta^{(j)}(-1),\qquad Z_j=\zeta(j),\qquad g=\gamma.
\]
Spectral derivatives are denoted by parentheses on zeta, whereas
subscripts on a multivariable analytic function mean partial derivatives.
The classical Mellin continuation \cite{BorweinDilcher} of
\[
 T(A,B,C)=\sum_{m,n\geq1}m^{-A}n^{-B}(m+n)^{-C}
\]
has the following local form, already established in the incoming
independent-order report \cite{RepoIndependent}:
\begin{equation}\label{qt:polar}
 T(A,B,C)=\zeta(A)\zeta(B)
 -\frac{C\zeta(B-1)}{A+C}
 -\frac{C\zeta(A-1)}{B+C}+C R(A,B,C),
\end{equation}
where $R$ is holomorphic near the origin and symmetric in $A,B$.
The two indicated polar quotients must be restricted before evaluating
at the origin. Thus for $(a+c)(b+c)\ne0$ the function
$W_{a,b,c}(t)=T(at,bt,ct)$ is holomorphic near zero, although $T$ is not
jointly holomorphic there.

\subsection{Six cubic coefficients and three exact constraints}

All derivatives in the next display are evaluated at $(0,0,0)$:
\[
 \alpha=R_{AAA},\quad \beta=R_{AAB},\quad
 \delta=R_{AAC},\quad \varepsilon=R_{ABC},\quad
 \varphi=R_{ACC},\quad \chi=R_{CCC}.
\]
Symmetry gives the corresponding derivatives with $A,B$ interchanged.
These six quantities describe the cubic homogeneous part of $R$.

\begin{proposition}\label{qt:constraints}
One has
\begin{equation}\label{qt:constraints-eq}
 \chi=\frac{e_4-d_4}{4},\qquad
 \delta=\frac{d_2^2-d_4}{4},\qquad
 \alpha+3\varphi=-\frac{Ld_3}{2}-d_4.
\end{equation}
If $\Omega_4=W_{1,1,1}^{(4)}(0)$, then
\begin{equation}\label{qt:diagonal}
 \boxed{\Omega_4=24(\beta+\varepsilon)-8Ld_3+12d_2^2-16d_4.}
\end{equation}
\end{proposition}
\begin{proof}
The exact axis $T(0,0,C)=\zeta(C-1)-\zeta(C)$ gives
\[
 C R(0,0,C)=\zeta(C-1)-\zeta(C)-\frac5{12}.
\]
Its coefficient of $C^4$ proves the first assertion.
The Euler stuffle identity on the plane $B=0$ gives
\begin{align}\label{qt:stuffle}
 C R(A,0,C)+A R(C,0,A)
 ={}&\zeta(A)\zeta(C)-\zeta(A+C)
 +\tfrac12\bigl(\zeta(A)+\zeta(C)\bigr)\notag\\
 &+\zeta(-1)+\zeta(A-1)+\zeta(C-1).
\end{align}
The coefficient of $A^2C^2$ is $\delta$ on the left and
$(d_2^2-d_4)/4$ on the right. The coefficient of $A^3C$ is
$\alpha/6+\varphi/2$ on the left and
$d_1d_3/6-d_4/6$ on the right. Since $d_1=-L/2$, this proves
the remaining two constraints.

On the diagonal, \eqref{qt:polar} becomes
$W_{1,1,1}(t)=\zeta(t)^2-\zeta(t-1)+tR(t,t,t)$.
Its fourth derivative is
\[
 -d_4-4Ld_3+6d_2^2-e_4
 +8\alpha+24\beta+24\delta+24\varepsilon+24\varphi+4\chi.
\]
Substituting \eqref{qt:constraints-eq} proves \eqref{qt:diagonal}.
\end{proof}

The three constraints have rank three in the six-dimensional space of
formal cubic jets symmetric in $A,B$. Thus that formal reconstruction
problem has exactly three free coordinates. This dimension statement
does not imply arithmetic independence of the evaluated constants.

\subsection{An explicit formula for every fourth ray derivative}

\begin{theorem}\label{qt:main}
For $a,b,c\in\mathbb C$ with $(a+c)(b+c)\ne0$,
\begin{align}\label{qt:all-rays}
 W_{a,b,c}^{(4)}(0)={}&abc^2\Omega_4
 +4c(a+b)(a^2-ab+b^2-c^2)\alpha\notag\\
 &+12abc(a+b-2c)\beta\notag\\
 &-2\bigl[ab(a^2+b^2)-4abc^2+c^3(a+b)\bigr]Ld_3\notag\\
 &+3\bigl[2a^2b^2+c^2(a^2+b^2-4ab)\bigr]d_2^2\notag\\
 &+\bigl[-\tfrac12(a^4+b^4)-3c^2(a^2+b^2)
              +16abc^2-4c^3(a+b)-c^4\bigr]d_4\notag\\
 &+\left[c^4-\frac{cb^4}{a+c}-\frac{ca^4}{b+c}\right]e_4.
\end{align}
The coordinates $\alpha,\beta,\Omega_4$ have the convergent
Gamma and polylogarithmic representations below; no unspecified partial
derivative is required to evaluate the formula.
\end{theorem}
\begin{proof}
Expanding the holomorphic restriction of \eqref{qt:polar} gives
\begin{align*}
 W_{a,b,c}^{(4)}(0)={}&-\tfrac12(a^4+b^4)d_4
 -2ab(a^2+b^2)Ld_3+6a^2b^2d_2^2\\
 &-c\left(\frac{b^4}{a+c}+\frac{a^4}{b+c}\right)e_4
 +4c\alpha(a^3+b^3)+12c\beta ab(a+b)\\
 &+12c^2\delta(a^2+b^2)+24abc^2\varepsilon
 +12c^3\varphi(a+b)+4c^4\chi.
\end{align*}
Eliminate $\delta,\varphi,\chi$ using \eqref{qt:constraints-eq},
then eliminate $\varepsilon$ with \eqref{qt:diagonal}. Factoring the
$\alpha$ coefficient proves the claimed formula.
\end{proof}

For instance,
\begin{align}
 W_{1,1,2}^{(4)}(0)={}&4\Omega_4-48\alpha-48\beta-4Ld_3
                      -18d_2^2-41d_4+\tfrac{44}{3}e_4,
 \label{qt:112}\\
 W_{1,2,2}^{(4)}(0)={}&8\Omega_4-24\alpha-48\beta-4Ld_3
                      -12d_2^2-\tfrac{105}{2}d_4+\tfrac{29}{6}e_4.
 \label{qt:122}
\end{align}
The three positive anchor rays $(1,1,1),(1,1,2),(1,2,2)$ recover
all three formal coordinates: the determinant of their coefficient
matrix in $(\Omega_4,\alpha,\beta)$ is $1152$.

\subsection{Convergent representatives: a Gamma moment and two kernels}

The incoming boundary-jet hierarchy \cite{RepoIndependent} already gives, with
\[
 \rho(n)=\log\Gamma(n+1)-(n+\tfrac12)\log n+n-\frac L2-\frac1{12n},
\]
the identity
\begin{equation}\label{qt:alpha}
 \alpha=-\sum_{n\ge1}\rho(n)\log^3n-e_4-\tfrac12d_4-e_3-\frac{\gamma_3}{12}.
\end{equation}
The series converges absolutely because $\rho(n)=O(n^{-3})$.
Its use as a quartic coordinate is new here; the all-order boundary
identity itself is credited to the incoming report.

For the other two coordinates, define for $x>0$
\[
 f_j(x)=\left.\partial_s^j\operatorname{Li}_s(e^{-x})\right|_{s=0},
 \qquad
 p_1(x)=\frac{g+\log x}{x},\quad
 p_2(x)=\frac{(g+\log x)^2+Z_2}{x}.
\]
Put
\begin{align*}
 S_{21}(x)&=p_2(x)p_1(x)+p_2(x)d_1+p_1(x)d_2
                  -x\bigl(p_2(x)e_1+p_1(x)e_2\bigr)+d_2d_1,\\
 S_{11}(x)&=p_1(x)^2+2p_1(x)d_1-2xp_1(x)e_1+d_1^2.
\end{align*}
The ordinary convergent integrals are
\begin{align}\label{qt:J21}
 J_{21}={}&\int_0^1\frac{f_2(x)f_1(x)-S_{21}(x)}x\,dx
                  +\int_1^\infty\frac{f_2(x)f_1(x)}x\,dx,\\
 K_{11}={}&\int_0^1\frac{\log x+g}{x}
                    \bigl(f_1(x)^2-S_{11}(x)\bigr)\,dx
                  +\int_1^\infty\frac{\log x+g}{x}f_1(x)^2\,dx.
 \label{qt:K11}
\end{align}
Set $\mu_2=g^2+Z_2$ and $\mu_3=g^3+3gZ_2+2Z_3$, and define
\begin{align}\label{qt:constants}
 E_{21}={}&gd_1d_2-(\mu_2+2g+2)d_1-(g+1)d_2
                -\frac{\mu_3e_1}{3}-\frac{\mu_2e_2}{2}\notag\\
 &-\frac38-\frac{3g}{4}-\frac{\mu_2}{4}
                -\frac{g^2}{2}-\frac{g\mu_2}{2},\\
 E_{111}={}&-\frac{g^3}{2}-\frac{3g^2}{4}-\frac{3g}{4}-\frac38
               -2(g^2+2g+2)d_1\notag\\
 &+\frac{g^2-Z_2}{2}d_1^2-\frac{2(g^3-Z_3)}{3}e_1.
\end{align}

\begin{proposition}\label{qt:integrals}
The integral and constant formulas give
\begin{equation}\label{qt:beta-epsilon}
 \boxed{\beta=J_{21}+E_{21},\qquad
        \varepsilon=K_{11}+E_{111}.}
\end{equation}
Together with \eqref{qt:diagonal}, these evaluate $\Omega_4$ using
two ordinary, absolutely convergent one-dimensional integrals.
\end{proposition}
\begin{proof}
For parameters near zero put $G(C)=1/\Gamma(1+C)$ and
\begin{align*}
 S(A,B;x)={}&\Gamma(1-A)\Gamma(1-B)x^{A+B-2}
 +\Gamma(1-A)\zeta(B)x^{A-1}
 +\Gamma(1-B)\zeta(A)x^{B-1}\\
 &-\Gamma(1-A)\zeta(B-1)x^A
 -\Gamma(1-B)\zeta(A-1)x^B+\zeta(A)\zeta(B),\\
 I(A,B,C)={}&\int_0^1x^{C-1}
  \bigl(\operatorname{Li}_A(e^{-x})\operatorname{Li}_B(e^{-x})-S(A,B;x)\bigr)dx\\
 &+\int_1^\infty x^{C-1}
                    \operatorname{Li}_A(e^{-x})\operatorname{Li}_B(e^{-x})dx,\\
 X(A,B,C)={}&\frac{\Gamma(1-A)\Gamma(1-B)}{A+B+C-2}
  +\frac{\Gamma(1-A)\zeta(B)}{A+C-1}
  +\frac{\Gamma(1-B)\zeta(A)}{B+C-1}.
\end{align*}
Jonqui\`ere's expansion at zero \cite[Eq.~25.12.12]{DLMFPolylog} shows that the subtracted integrand
and each fixed parameter derivative are dominated by an integrable
power times a polynomial in $|\log x|$ on a sufficiently small common
polydisc. The tail is exponentially decreasing. Thus $I$ is holomorphic.
Integrating the six subtractions explicitly in the classical Mellin
representation, then using \eqref{qt:polar}, gives
\begin{align}\label{qt:R-integral}
 R(A,B,C)={}&G(C)\bigl(I(A,B,C)+X(A,B,C)\bigr)
            +\frac{G(C)-1}{C}\zeta(A)\zeta(B)\notag\\
 &-\frac{\Gamma(1-A)G(C)-1}{A+C}\zeta(B-1)
  -\frac{\Gamma(1-B)G(C)-1}{B+C}\zeta(A-1).
\end{align}
Every displayed quotient is removable at the origin. For example,
$\Gamma(1-A)G(-A)=1$ proves divisibility by $A+C$.

The derivative $I_{AAB}(0)$ is exactly $J_{21}$, since differentiating
the six subtractions gives $S_{21}$. Likewise
$I_{ABC}(0)+gI_{AB}(0)=K_{11}$, since $G'(0)=g$ and the mixed
subtraction is $S_{11}$. The remaining derivatives in
\eqref{qt:R-integral} are finite Gamma and zeta algebra. In particular,
\[
 \left.\partial_A\partial_C
       \frac{\Gamma(1-A)G(C)-1}{A+C}\right|_0
       =\frac{g^3-Z_3}{3}.
\]
For $C=0$, the derivatives of
$(\Gamma(1-A)-1)/A$ of orders $1,2$ are $\mu_2/2,\mu_3/3$.
These evaluations and the elementary denominators in $X$ give exactly
$E_{21}$ and $E_{111}$.

Finally $f_2f_1-S_{21}=O(x(1+|\log x|^2))$ and
$f_1^2-S_{11}=O(x(1+|\log x|))$ near zero. Their displayed
integrals therefore converge absolutely in the ordinary sense.
\end{proof}

\subsection{Exact cancellation families}

\begin{corollary}[Cyclic quartic reduction]\label{qt:cyclic}
Suppose all pairwise sums of $a,b,c$ are nonzero, and let
$s_1=a+b+c$, $s_2=ab+bc+ca$, $s_3=abc$. Then
\begin{align}\label{qt:cyclic-eq}
 \sum_{\mathrm{cyc}} W_{a,b,c}^{(4)}(0)={}&s_1s_3\Omega_4
 +(-4s_1^2s_2+8s_2^2+12s_1s_3)Ld_3\notag\\
 &+(12s_2^2-36s_1s_3)d_2^2\notag\\
 &+(-2s_1^4+4s_1^2s_2-2s_2^2+24s_1s_3)d_4.
\end{align}
Both off-diagonal coordinates and $e_4$ cancel identically.
\end{corollary}
\begin{proof}
The $\alpha$ coefficient is
$4\sum_{\rm cyc}c(a^3+b^3)-4\sum_{\rm cyc}c^3(a+b)=0$.
The $\beta$ coefficient is $12abc\sum_{\rm cyc}(a+b-2c)=0$.
For the $e_4$ coefficient, the terms with denominator $a+c$
combine to $-b^4$ and cancel the sum of fourth powers. The remaining
symmetric polynomials yield \eqref{qt:cyclic-eq}.
\end{proof}

There is also a useful one-parameter source of completely evaluated
differences. On the additive plane $c=a+b$, the nonclassical part of
\eqref{qt:all-rays} is
\[
 ab(a+b)^2(\Omega_4-12\alpha-12\beta).
\]
Any two admissible rays on this plane therefore have a weighted
difference involving only ordinary zeta derivatives. A small example is
\begin{equation}\label{qt:short-difference}
 \boxed{W_{1,2,3}^{(4)}(0)-\frac92 W_{1,1,2}^{(4)}(0)
 =-20Ld_3+24d_2^2-76d_4+\frac{12}{5}e_4.}
\end{equation}
This is an evaluated identity with no undetermined Tornheim coordinate.

For a single ray with $abc\ne0$, the two displayed off-diagonal
coefficients vanish simultaneously if and only if $a=b=c$.
Indeed the $\beta$ coefficient forces $a+b=2c$; substituting this
into the $\alpha$ coefficient gives a nonzero multiple of
$(a-b)^2$. Thus the nontrivial cubic cancellation conic from the
incoming report does not persist as a conic at the fourth layer.
This assertion concerns the stated formal-coordinate elimination.
An additional arithmetic relation among the constants could enlarge
the family of values reducible in a chosen constant algebra.

\subsection{All-order jet structure and the limit of cyclic diagonal reduction}

The quartic formula is the last layer at which the axis and Euler
stuffle constraints force a cyclic sum to depend on only one diagonal
coordinate. The following exact polynomial statement explains why.

Fix an integer $n\ge3$. Subtract two possible degree-$n$ homogeneous
jets of \eqref{qt:polar} whose axis and Euler boundary restrictions
agree. Their difference is a homogeneous polynomial $U(A,B,C)$
of degree $n$ satisfying
\begin{equation}\label{qt:residual-conditions}
 U(A,B,C)=U(B,A,C),\qquad C\mid U,\qquad
 U(A,0,C)+U(C,0,A)=0.
\end{equation}
These are conditions on a formal reconstruction space, not additional
functional equations asserted for arbitrary perturbations of $T$.

\begin{theorem}[Normal form of the residual jet]\label{qt:all-order-space}
Every polynomial satisfying \eqref{qt:residual-conditions} has a unique
representation
\begin{equation}\label{qt:residual-form}
 U=C\bigl[A(A-C)V(A,C)+B(B-C)V(B,C)\bigr]+ABC\,Q(A,B,C),
\end{equation}
where $V$ is a symmetric two-variable homogeneous polynomial of degree
$n-3$, and $Q$ is homogeneous of degree $n-3$ and symmetric in $A,B$.
Consequently this residual space has dimension
\begin{equation}\label{qt:jet-dimension}
 D_n=\left\lfloor\frac{n^2-1}{4}\right\rfloor.
\end{equation}
For $n=1,2$ the corresponding residual space is zero.
\end{theorem}
\begin{proof}
The boundary polynomial $F(A,C)=U(A,0,C)$ is antisymmetric in $A,C$.
It is divisible by $C$, hence by $A$ as well, and antisymmetry also
implies divisibility by $A-C$. Therefore
$F(A,C)=AC(A-C)V(A,C)$, where $V$ is symmetric and has degree
$n-3$. This uniquely determines $V$. Subtract the first term of
\eqref{qt:residual-form} from $U$. The difference vanishes on both
$A=0$ and $B=0$, by symmetry, and remains divisible by $C$.
Hence it is uniquely $ABCQ$, with the required symmetry.
Conversely the displayed normal form satisfies all three conditions.

The symmetric two-variable polynomials of degree $n-3$ have dimension
$\lfloor(n-1)/2\rfloor$. Polynomials in three variables of that degree
which are symmetric in the first two have dimension
$\lfloor(n-1)^2/4\rfloor$, by counting the exponent pairs of $A,B$.
Their sum is \eqref{qt:jet-dimension}. If $n<3$, the two divisibility
arguments force $U=0$ directly.
\end{proof}

\begin{corollary}[Cyclic residual dimension]\label{qt:cyclic-dimension}
The cyclic symmetrization of the residual space is precisely
\[
 \left\{ABC\,S(A,B,C): S\text{ is symmetric in all three variables
 and homogeneous of degree }n-3\right\}.
\]
Its dimension is
\begin{equation}\label{qt:cyclic-count}
 p_{\le3}(n-3)=\left\lfloor\frac{n^2+3}{12}\right\rfloor,
\end{equation}
where $p_{\le3}(m)$ is the number of partitions of $m$ into at most
three parts. In particular this dimension is one at orders three and
four, and is two at order five.
\end{corollary}
\begin{proof}
The terms involving $V$ in \eqref{qt:residual-form} cancel in pairs
under cyclic summation, because $V(A,C)=V(C,A)$. The remaining sum is
$ABC\sum_{\rm cyc}Q(A,B,C)$. Since $Q$ is symmetric in its first
two variables, this sum is fully symmetric. Conversely any symmetric
$S$ occurs by taking $Q=S/3$ and $V=0$. Symmetric homogeneous
polynomials of degree $m$ have basis $s_1^i s_2^j s_3^k$ with
$i+2j+3k=m$. Counting these triples, or expanding
$(1-t)^{-1}(1-t^2)^{-1}(1-t^3)^{-1}$, gives
$p_{\le3}(m)=\lfloor(m^2+6m+12)/12\rfloor$, proving
\eqref{qt:cyclic-count}.
\end{proof}

Evaluation on the diagonal is a nonzero linear functional on this
cyclic residual space. It therefore determines the whole space when
the dimension is one, explaining the complete cyclic formulas at
orders three and four. The following explicit fifth-order polynomial
shows what fails next:
\begin{equation}\label{qt:fifth-obstruction}
 U_5(A,B,C)=ABC\bigl[(A+B+C)^2-3(AB+BC+CA)\bigr].
\end{equation}
It satisfies the boundary conditions, vanishes on the entire diagonal,
and has nonzero cyclic sum $3U_5$. Consequently the axis, symmetry,
Euler stuffle, and diagonal data alone cannot force a fifth-order
cyclic formula involving just the diagonal fifth derivative and
ordinary zeta jets. This is a precise limitation of those input
identities; other functional equations may provide additional
constraints.

\subsection{Verification and remaining arithmetic questions}

The companion implementation \texttt{check\_quartic.py} checks
the full rational-function identity, cyclic elimination, diagonal and
Euler-diagonal cases, and both elementary integral correction terms.
It evaluates $\alpha$ with a finite Stirling expansion and computes
$J_{21},K_{11}$ by an integrated Jonqui\`ere series on $(0,1)$ plus
exponentially convergent series on $(1,\infty)$. A separate evaluator
computes the fourth ray derivative by a fifteen-node symmetric
difference of the original Mellin continuation. These routes give
\begin{align*}
 \alpha&=13.1294216474884019236329146797203\ldots,\\
 \beta&=6.90381599352243710775587231649323\ldots,\\
 \varepsilon&=4.68913245627221585695112599939183\ldots,\\
 \Omega_4&=798.777376093189588968747119645501\ldots.
\end{align*}
At 55 decimal digits of working precision, four independently evaluated
rays $(1,1,1)$, $(1,0,2)$, $(1,2,2)$, and $(1,-2,3)$ had absolute discrepancies
below $8.2\cdot10^{-34}$. These decimal checks use mpmath and are not
interval certificates; the exact identities rest on the proofs above.

The result completes the quartic finite-coordinate problem posed in the
incoming reports. It does not provide a finite expression for
$\Omega_4,\alpha,\beta$ in a smaller algebra of ordinary constants.
Natural next questions are: relate the two new integral coordinates to
higher harmonic Stieltjes coefficients; identify the dimension and
canonical anchor directions at every order; obtain general cyclic
elimination laws; and formalize the finite jet algebra separately from
the Mellin continuation and normal-convergence lemmas.

